\section{Stable fields along Gaussian observations}
\label{sec:stability}

The observation argument requires control of every approximate solution
of the field equation at the fields visited by the observation process.
We first prove this stability statement, including perturbations of the
diagonal variance matrix, and then derive a deterministic existence,
uniqueness, and inverse estimate for the field equation.

For a symmetric matrix $J$, a field $h_0\in\R^n$, and $y\in\R^n$, define
\begin{equation}
 \begin{split}
 m(y)&=\tanh y,\qquad q(y)=\av{m(y)^2},\qquad
 b(y)=1-q(y),\\
 V_y&=D_{1-m(y)^2},\qquad B(y)=jb(y),\\
 F_{J,h_0}(y)&=y-h_0+(B(y)I-J)m(y).
 \end{split}
 \label{eq:field-map}
\end{equation}
For fixed $J$, abbreviate $F_h=F_{J,h}$; we also suppress $h$ when
there is no ambiguity. Hyperbolic functions act coordinatewise. To measure
stability while allowing the diagonal factor to vary, define, for each
nonnegative diagonal matrix $A$,
\begin{equation}
 \mathsf H_J(y,A)
 =I+A^{1/2}\left(B(y)I-J-\frac{2j}{n}m(y)m(y)^{\mathsf T}\right)A^{1/2}.
 \label{eq:field-hessian}
\end{equation}
The matrix $\mathsf H_J(y,V_y)$ is the symmetric version of the
possibly nonsymmetric derivative $F'_{J,h_0}(y)$. The external field
$h_0$ enters the residual $F_{J,h_0}(y)$ but not
$\mathsf H_J(y,A)$. This separation will let us use the same matrix
witness when the observation field moves.

\begin{proposition}[Stability along observations]
\label{prop:stable-fields}
Fix $j\in(0,1)$ and $0<T<\infty$.  There are constants
$\eta,c,\epsilon_A,\rho_0,a>0$, depending only on $j,T$, and events
$\mathcal J_n$ for the zero-diagonal Gaussian matrix $J$, such that
$\Prob(J\in\mathcal J_n)\longrightarrow1$ and the following assertion
holds for every $J\in\mathcal J_n$.  Let $X\sim\mu_0$, let $B_0$ be an
independent standard $n$-dimensional Brownian motion, and put
$Y(t)=tX+B_0(t)$.  With conditional probability at least $1-e^{-an}$,
simultaneously for $0\le t\le T$, $y\in\R^n$, and all diagonal $A$,
\begin{equation}
 \left.
 \begin{gathered}
 \norm{F_{J,Y(t)}(y)}\le2\rho_0\sqrt n,\\
 0\le A\le(1+\eta)I,\qquad
 \norm{A-V_y}_{\Frob}\le\epsilon_A\sqrt n
 \end{gathered}
 \right\}
 \quad\Longrightarrow\quad
 \mathsf H_J(y,A)\succ cI.
 \label{eq:stable-field}
\end{equation}
We may require $j(1+\eta)^2<1$ and
$\norm J\le2\sqrt j+\theta$ on $\mathcal J_n$, for any sufficiently
small fixed $\theta>0$ chosen in the proof.
\end{proposition}

\begin{corollary}[Stopping formulation]\label{cor:stopping}
For every $J\in\mathcal J_n$ there is an observation stopping time $\tau_0$
such that
\[
 \Prob(\tau_0\le T\mid J)\le e^{-an},
\]
and, for every $t\in[0,T]$ with $t<\tau_0$, the implication \eqref{eq:stable-field} holds
at $Y(t)$ with residual threshold $\rho_0\sqrt n$ in place of
$2\rho_0\sqrt n$. The stopping rule depends on $J$ and the observations,
and is independent of any test function.
\end{corollary}

\begin{proof}
We express failure of stability through a witness set that is fixed
before observing the path. Let $\mathcal W_J$ consist of pairs $(y,A)$ satisfying
$0\le A\le(1+\eta)I$ and $\norm{A-V_y}_{\Frob}\le\epsilon_A\sqrt n$
for which $\mathsf H_J(y,A)\not\succ cI$. Define the least
normalized residual among these witnesses by
\[
 d_J(h)=\inf_{(y,A)\in\mathcal W_J}
             \frac{\norm{F_{J,h}(y)}}{\sqrt n},
 \qquad \inf\varnothing=+\infty.
\]
If $\mathcal W_J$ is nonempty, $d_J$ is finite and $n^{-1/2}$-Lipschitz,
because $F_{J,h}(y)=F_{J,0}(y)-h$. Let
\[
 \tau_0=\inf\{t\in[0,T]:d_J(Y(t))\le\rho_0\},
\]
with value $+\infty$ when the set is empty. Continuity makes the first
hitting time an observation stopping time.
Before the hit, every violating witness has residual greater than
$\rho_0\sqrt n$. At a hit, the definition of the infimum supplies a
violating witness with residual strictly below $2\rho_0\sqrt n$;
attainment of the infimum is unnecessary.
Proposition~\ref{prop:stable-fields} therefore bounds the hitting probability.
\end{proof}

The proof of Proposition~\ref{prop:stable-fields} has three parts. A
Gaussian integral bounds the weight of unstable configurations; its
scalar rate is controlled by the entropy inequality, and its conditional
matrix factor is controlled near the equality laws. A local-volume
estimate then shows that any unstable approximate zero would contribute
too much to that integral. Finally we transfer the resulting estimate to
the original observation law, uniformly over the time interval.

\subsection{A Gaussian integral and its scalar rate}

We carry out the integral calculation first under a planted Gaussian
law, whose relation to the original law will be proved when we assemble
the path estimate. Under this law,
\begin{equation}
 \wt J=W+\frac jn\one\one^{\mathsf T},\qquad
 h_0=t\one+\sqrt t\,g,
 \label{eq:planted-field}
\end{equation}
where $W$ is centered GOE of scale $j$ and $g$ is independent standard
Gaussian. We retain the diagonal of $\wt J$ during this calculation.
For a Borel set $\mathcal E\subseteq\R^n\times\mathrm{Sym}_n$, define
\begin{equation}
 I_{t,\sigma}(\mathcal E)
 =(2\pi\sigma^2)^{-n/2}\int
 \one_{\{(y,\wt J)\in\mathcal E\}}
 \exp\left\{-\frac{\norm{F_{\wt J,h_0}(y)}^2}{2\sigma^2}
              +\frac{njb(y)^2}{2}\right\}\,\dd y,
 \qquad \sigma>0.
 \label{eq:field-integral}
\end{equation}
The weight $e^{njb(y)^2/2}$ compensates for the exponential scale of
the local determinant: Lemma~\ref{lem:local-field-volume} will use the
regularized determinant bound to turn an approximate zero into a lower
bound for this integral. The sets to which we apply the integral depend
on $y,\wt J$ and are independent of $h_0$. Thus the observation can be
averaged while retaining the restriction to $\mathcal E$.

For fixed $y,t,\sigma$, introduce the scalar quantities
\begin{equation}
 \begin{gathered}
 m=\tanh y,\quad q=\av{m^2},\quad b=1-q,\quad B=jb,
 \quad M=\av m,\\
 d=t+jM,\qquad s=t+\sigma^2+jq,\qquad S=s+jq,\\
 a_y=\av{(y-d\one)m},\qquad z=y-d\one+Bm,\qquad
 \psi=\frac{j(a_y-sb)^2}{2sS}.
 \label{eq:scalar-parameters}
 \end{gathered}
\end{equation}
Write $\varphi_{d,s}$ for the density of $N(d,s)$, with $s$ denoting
the variance, and define
\begin{equation}
 G_{t,\sigma,n}(y)
 =\prod_{i=1}^n\varphi_{d,s}(y_i)\,e^{n\psi}.
 \label{eq:scalar-density}
\end{equation}
The Gaussian parameters in this expression depend on $y$ through the
preceding averages; the product is therefore not a fixed product density.

\begin{lemma}[Conditional Gaussian reduction]
\label{lem:gaussian-reduction}
With the preceding definitions, and with $g'$ an independent standard
Gaussian vector,
\begin{equation}
 \E I_{t,\sigma}(\mathcal E)
 =\int \sqrt{\frac sS}\,G_{t,\sigma,n}(y)
 \Prob\left((y,\wt J)\in\mathcal E\,\middle|\,
       Wm+\sqrt{t+\sigma^2}\,g'=z\right)\,\dd y.
 \label{eq:conditional-integral}
\end{equation}
The conditional distribution in this identity is the usual Gaussian
conditional distribution and is defined for every $y$, since
$\sigma>0$.
\end{lemma}

\begin{proof}
We first evaluate the integral without its indicator. For fixed $m$,
the GOE covariance is
\[
 \Cov(Wm)=jqI+\frac jnmm^{\mathsf T}.
\]
Convolution with the heat kernel of variance $\sigma^2$ and with
$\sqrt t\,g$ therefore produces the Gaussian vector
$Wm+\sqrt{t+\sigma^2}\,g'$, of covariance
$sI+jmm^{\mathsf T}/n$.  Its determinant is $s^nS/s$, and its inverse
is
\[
 \frac1sI-\frac{j}{nsS}mm^{\mathsf T}.
\]
Moreover,
$F_{\wt J,h_0}(y)=z-Wm-\sqrt t\,g$.  The Gaussian density at $z$,
multiplied by $e^{njb^2/2}$, differs from
$\prod_i\varphi_{d,s}(y_i)$ by the prefactor $\sqrt{s/S}$ and the
exponential whose logarithm divided by $n$ is
\[
 -\frac{Ba_y}{s}-\frac{B^2q}{2s}
 +\frac{j(a_y+Bq)^2}{2sS}+\frac{jb^2}{2}
 =\frac{j(a_y-sb)^2}{2sS}.
\]
This calculation gives the identity for the unrestricted integral.
To retain the indicator, condition the joint Gaussian pair formed by
$W$ and the convolved observation on the value $z$ of the latter.
The resulting conditional probability is exactly the factor in
\eqref{eq:conditional-integral}. All integrands are nonnegative, so
Tonelli's theorem justifies the interchange of integrations.
\end{proof}

The scalar entropy inequality controls the remaining integral in
\eqref{eq:conditional-integral}. To apply it, replace empirical averages
in \eqref{eq:scalar-parameters} by integration against a probability
measure $P$ with finite second moment, keeping the notation
$M,q,b,d,s,S,a_y,\psi$. Where needed, write $a(P)$ instead of $a_y$.
Lemma~\ref{lem:scalar-entropy} then gives
\begin{equation}
 \psi(P,t,\sigma)-D(P\Vert N(d,s))\le0
 \qquad (s>0).
 \label{eq:rate-nonpositive}
\end{equation}
At $\sigma=0$, Lemma~\ref{lem:gaussian-equality} identifies all equality
cases: they form the curve
\[
 P_t=N(s_*(t),s_*(t)),\qquad
 s_*(t)=t+j\E_{P_t}\tanh y,
\]
with the point-mass convention $P_0=\delta_0$.  In particular,
$s_*(t)\le t/(1-j)$ and
$q(P_t)=(s_*(t)-t)/j\le t/(1-j)$.

\begin{lemma}[Uniform scalar rate bounds]
\label{lem:scalar-rates}
Fix $j\in(0,1)$ and $T<\infty$.
\begin{enumerate}[label=\textup{(\roman*)}]
\item For every $s_{\min}>0$ and $L<\infty$, one can choose $R<\infty$
such that, for all sufficiently large $n$, uniformly over
$0\le t\le T$ and $0<\sigma\le1$,
\[
 \int_{\{s\ge s_{\min},\,\av{y^2}>R^2\}}
 G_{t,\sigma,n}(y)\,\dd y\le e^{-Ln}.
\]
For every $\varepsilon>0$, the unrestricted integral over
$\{s\ge s_{\min}\}$ is at most $e^{\varepsilon n}$ for all
sufficiently large $n$.
\item Fix $t_0>0$.  Outside any prescribed open neighborhood of the
equality curve $\{(P_t,t):t_0\le t\le T\}$, there are
$a,\sigma_*>0$ such that the integral of $G_{t,\sigma,n}$ is at most
$e^{-an}$, uniformly over $t_0\le t\le T$ and
$0<\sigma\le\sigma_*$, for all sufficiently large $n$.
The neighborhood can impose weak closeness of the empirical law and
closeness of any finite list of continuous, at most linearly growing
averages, including all parameters in
\eqref{eq:scalar-parameters}.
\item For every $q_0>0$, there are $t_0,\sigma_*,a>0$ such that
\[
 \int_{\{q\ge q_0\}}G_{t,\sigma,n}(y)\,\dd y\le e^{-an}
\]
uniformly over $0\le t\le t_0$ and $0<\sigma\le\sigma_*$, for all
sufficiently large $n$.
\item Fix $t_0>0$ and $\nu>0$.  For all sufficiently large $R_1$ there
is $a>0$ such that
\[
 \int_{\{\av{y^2\one_{\{|y|>R_1\}}}>\nu\}}
 G_{t,\sigma,n}(y)\,\dd y\le e^{-an},
\]
uniformly over $t_0\le t\le T$ and $0<\sigma\le1$, for all
sufficiently large $n$.
\end{enumerate}
All constants are independent of $n,t,\sigma$ in the indicated ranges.
In \textup{(ii)}, neighborhoods are interpreted on each fixed
second-moment ball; the complement of a sufficiently large such ball
is covered by \textup{(i)}.
\end{lemma}

\begin{proof}
We first obtain a Gaussian tail envelope, then prove a compact-set
bound by a finite cover. Strict negativity away from the equality laws
will give the second and third assertions; an exponential tilt will give
the fourth.

Throughout, $|d|\le T+j$ and $s\le T+1+j$. Also
$a(P)^2\le q\int(y-d)^2\,\dd P$ and $jq/S\le1/2$.  The elementary
inequality $(r-u)^2\le\frac32r^2+3u^2$ consequently gives
\[
 n\psi\le\frac{3jq}{4sS}\sum_i(y_i-d)^2+\frac{3nj}{2}
 \le\frac{3}{8s}\sum_i(y_i-d)^2+\frac{3nj}{2}.
\]
On $s\ge s_{\min}$, it follows that
\begin{equation}
 G_{t,\sigma,n}(y)\le
 \exp\left(Cn-c\sum_i y_i^2\right),
 \label{eq:scalar-tail-envelope}
\end{equation}
where $c>0$ and $C<\infty$ depend only on $j,T,s_{\min}$.
On $\sum_i y_i^2>nR^2$, use half of the remaining Gaussian square
for integration and the other half for the required exponential decay.
This proves the first assertion of \textup{(i)}.

The estimates on bounded second moments follow from one compactness
argument, which we now give in full. Put
\[
 \mathcal P_R=\left\{P:\int y^2\,\dd P\le R^2\right\}.
\]
Tightness and preservation of the second-moment bound under weak limits
make this set weakly compact. Integration of a continuous function of at
most linear growth is continuous on it: the second-moment bound makes
the contribution of $|y|>K$ uniformly $O(K^{-1})$, while bounded
truncations are weakly continuous. In particular, $d,s,a(P),\psi$ are
continuous on
$\mathcal P_R\times[0,T]\times[0,1]$ where $s\ge s_{\min}$.

For every compact subset $\mathcal C$ of that region,
\begin{equation}
 \limsup_{n\to\infty}\frac1n\log
 \sup_{t,\sigma}\int_{\{(P_y,t,\sigma)\in\mathcal C\}}
 G_{t,\sigma,n}(y)\,\dd y
 \le\sup_{\mathcal C}\bigl\{\psi-D(P\Vert N(d,s))\bigr\},
 \label{eq:compact-rate-bound}
\end{equation}
where $P_y=n^{-1}\sum_i\delta_{y_i}$; empty integrals are zero.
To establish this bound, fix a strict finite upper bound $r$ for its
right-hand side. At each $\xi=(P,t,\sigma)\in\mathcal C$, choose a
bounded continuous test function $\varphi_\xi$ from the variational
formula for relative entropy so that
\[
 \psi(\xi)-\int\varphi_\xi\,\dd P
       +\log\int e^{\varphi_\xi}\,\dd N(d,s)<r.
\]
Such a test also exists when the entropy is infinite. Retain positive
slack in the inequality and then choose a neighborhood of $\xi$ on
which continuity controls $\psi$, the empirical average of
$\varphi_\xi$, and $d,s$.  Moreover, on
$\sum_i y_i^2\le nR^2$, the logarithmic density ratio between
$N(d',s')^{\otimes n}$ and the fixed law $N(d,s)^{\otimes n}$ is at
most $n$ times a quantity tending uniformly to zero as
$(d',s')\to(d,s)$; this follows directly by expanding their Gaussian
log densities.  For $y$ in that neighborhood, insert
$e^{\sum_i\varphi_\xi(y_i)}$ and bound its expectation under the
fixed product Gaussian.  The integral on the neighborhood is then at
most $e^{nr}$, after using less than the reserved slack for the three
continuity errors. Cover $\mathcal C$ by finitely many such
neighborhoods. An arbitrarily
small enlargement of $r$ absorbs their number and proves
\eqref{eq:compact-rate-bound}. Combining this bound with
\eqref{eq:rate-nonpositive} on the ball and
\eqref{eq:scalar-tail-envelope} outside it proves the second assertion
of \textup{(i)}.

To obtain a strictly negative rate, observe that the same variational
formula expresses relative entropy as a supremum of continuous
functions. Thus it is jointly lower semicontinuous in $P,d,s$ when
$s$ is bounded away from zero, and $\psi-D$ is upper semicontinuous.
At $\sigma=0$, its zero set is precisely the equality curve identified
above.  On a compact set disjoint from any open neighborhood of that
curve its supremum is strictly negative.  The same remains true for
all sufficiently small $\sigma$: otherwise a convergent sequence with
$\sigma\downarrow0$ and rates approaching zero would give an equality
point outside that neighborhood.  Apply
\eqref{eq:compact-rate-bound}, with $s_{\min}=t_0$, and then choose the
second-moment tail rate in \textup{(i)} strictly negative.  This proves
\textup{(ii)}.

Near time zero, the restriction on $q$ supplies the positive variance
needed for the same argument. For \textup{(iii)}, choose $t_0$ so small that
$t_0/(1-j)<q_0/2$.  The set $q\ge q_0$ contains no equality point at
$\sigma=0$ and $0\le t\le t_0$.  On this set $s\ge jq_0>0$, so the
preceding compactness argument and the tail envelope apply without
change, including at $t=0$.

It remains to control the contribution of large coordinates in
\textup{(iv)}. Fix $t_0,\nu>0$, and take
$0<\lambda<1/[8(T+1+j)]$.  Let
$h_{R_1}(y)=y^2\one_{\{|y|>R_1\}}$.  On a fixed second-moment ball,
repeat the proof of \eqref{eq:compact-rate-bound}, this time inserting
$e^{\lambda\sum_i h_{R_1}(y_i)}$.  Begin with any small positive rate
bound $\varepsilon$ in that proof and its finite collection of tests
$\varphi_\xi$.  Uniformly for $|d|\le T+j$ and
$t_0\le s\le T+1+j$,
\[
 \int\bigl(e^{\lambda h_{R_1}(y)}-1\bigr)\,\dd N(d,s)(y)
 \longrightarrow0\qquad(R_1\to\infty),
\]
by Gaussian tails; the choice of $\lambda$ leaves a uniformly negative
quadratic exponent.  The same holds with the integrand multiplied by
any of the finitely many bounded $e^{\varphi_\xi}$. The additional
logarithmic moment can therefore be made arbitrarily
small in every member of the finite cover. For sufficiently large $R_1$,
the tilted integral on the ball is at most $e^{2\varepsilon n}$.  On
$\av{h_{R_1}(y)}>\nu$, remove the tilt at cost
$e^{-\lambda\nu n}$.  Choosing $2\varepsilon<\lambda\nu/2$ gives a
strictly negative rate there.  The integral off the ball is estimated
without the tilt, using \textup{(i)} with a larger negative rate.
This proves \textup{(iv)}.
\end{proof}

\subsection{The conditional Hessian}

The scalar bounds concentrate the integral near the equality curve.
We now show that, there, the conditional Hessian is positive
with an exponential probability bound. The conclusion is simultaneous
over all admissible nearby diagonal matrices, which will be needed to
exclude every unstable approximate zero.

\begin{lemma}[Conditional stability near the scalar equality curve]
\label{lem:conditional-field-stability}
Fix $0<t_0\le T$, $R<\infty$, and $\eta>0$ with
$j(1+\eta)^2<1$.  There are $c_*,\epsilon_*,a_*,\sigma_*>0$,
a neighborhood $\mathcal U$ of
$\{(P_t,t):t_0\le t\le T\}$, a tail tolerance $\nu>0$, and a finite
$R_1$ such that the following holds.  Suppose
\[
 \av{y^2}\le R^2,\qquad (P_y,t)\in\mathcal U,\qquad
 \av{y^2\one_{\{|y|>R_1\}}}\le\nu,
 \qquad t_0\le t\le T,\quad0<\sigma\le\sigma_*.
\]
Under the Gaussian conditioning in
\eqref{eq:conditional-integral}, with probability at least
$1-e^{-a_*n}$, simultaneously for all diagonal $A$ satisfying
\[
 0\le A\le(1+\eta)I,\qquad
 \norm{A-V_y}_{\Frob}\le\epsilon_*\sqrt n,
\]
one has $\mathsf H_{\wt J}(y,A)\succeq c_*I$.
All constants are uniform over the displayed ranges and sufficiently
large $n$.  The tail cutoff $R_1$ may be required to be arbitrarily
large after choosing $\nu$; $\epsilon_*$ is then decreased if
necessary.
\end{lemma}

\begin{proof}
We first compute the Gaussian regression and the finite-rank correction
at an equality law. We then show that the empirical-law and diagonal
hypotheses preserve its positive margin, uniformly over all admissible
diagonals.

\paragraph{Gaussian regression and whitening.}
In a sufficiently small neighborhood of the equality curve, $s$ and
$q$ have positive lower bounds depending on $j,t_0,T$.
Indeed $s_*(t)\ge t_0$, and
$q(P_t)=\E_{P_t}\tanh^2 y$ is positive and continuous there.  Set
\[
 u=\frac{m}{\sqrt{nq}},\qquad
 Z_1=\frac{y-d\one}{\sqrt{ns}},\qquad P_u=I-uu^{\mathsf T}.
\]
The norms of $\one/\sqrt n,u,Z_1$ are bounded uniformly by the assumed
second-moment bound.  Write
\[
 W=P_uWP_u+uw^{\mathsf T}+wu^{\mathsf T}+\alpha uu^{\mathsf T}.
\]
The three blocks in this decomposition are independent before
conditioning, with $\Cov(w)=(j/n)P_u$ and $\Var(\alpha)=2j/n$.
The projection of the
observation in \eqref{eq:conditional-integral} perpendicular to $u$
observes $\sqrt{nq}\,w$ with independent noise of covariance
$(t+\sigma^2)P_u$; projecting along $u$ observes
$\sqrt{nq}\,\alpha$ with noise variance $t+\sigma^2$.  The elementary
Gaussian conditional-mean and conditional-variance formulas therefore
give the following representation, using a fresh centered GOE $W_0$:
\begin{equation}
 \begin{split}
 P_uWP_u&=P_uW_0P_u,\\
 w&=\bar w+\kappa P_uW_0u,\qquad
 \bar w=j\sqrt{q/s}\,Z_1-\frac{ja_y}{s}u,\qquad
 \kappa=\sqrt{\frac{t+\sigma^2}{s}},\\
 \alpha&=\frac{2j(a_y+Bq)}S
       +\sqrt{\frac{t+\sigma^2}S}\,u^{\mathsf T}W_0u.
 \end{split}
 \label{eq:conditional-W}
\end{equation}
For example, the perpendicular conditional mean is
$j\sqrt{q/s}\,P_uZ_1$, which equals the displayed $\bar w$ because
$u^{\mathsf T}Z_1=a_y/\sqrt{qs}$.  The parallel conditional mean follows
from $u^{\mathsf T}z=\sqrt{n/q}(a_y+Bq)$.  The parallel residual
variance is $(2j/n)(t+\sigma^2)/S$. These conditional means and
residual variances verify all three blocks
of \eqref{eq:conditional-W}.

For the moment, fix a deterministic admissible $A$. We compare the
conditional Hessian with the positive matrix built from the fresh GOE
by setting
\[
 \chi=\frac1n\Tr A,\qquad B_A=j\chi,\qquad
 S_A=I+B_AA-A^{1/2}W_0A^{1/2},\qquad
 \mathcal Y(p)=S_A^{-1/2}A^{1/2}p.
\]
Lemma~\ref{lem:matrix-path-stability} gives $S_A\succeq c_0I$ except
with exponentially small probability, where $c_0>0$ is uniform in
$A$.  Let $G_A=A^{1/2}S_A^{-1}A^{1/2}$.  Lemma~\ref{lem:bilinear},
applied to the three deterministic vectors
$\one/\sqrt n,u,Z_1$, says that their Gram entries after $\mathcal Y$
are as close as desired to their $A$-weighted Gram entries, and that
the extra vector $\mathcal Y(W_0u)$ has Gram entries as close as desired
to those of
\begin{equation}
 B_A\mathcal Y(u)+\sqrt{B_A}\,N,
 \label{eq:extra-whitened-vector}
\end{equation}
where $N$ is a unit vector orthogonal to the three deterministic
images in this comparison.  Specifically, its mixed entries are
$B_Ap^{\mathsf T}Au$ and its squared norm is
$B_A+B_A^2u^{\mathsf T}Au$. Here $N$ describes only Gram entries;
this notation does not assert a coupling with an additional random or
deterministic vector. The comparison concerns finite Gram matrices; the
associated symmetric finite-rank eigenvalues depend continuously on their
entries, as we verify when returning to the empirical parameters. Each fixed
Gram tolerance has failure probability at most $Ce^{-a n}$ uniformly
over the bounded deterministic vectors and over $A$.  Also
$u^{\mathsf T}W_0u$ is smaller than any fixed tolerance with
exponentially high probability, since its variance is $2j/n$.

\paragraph{The correction at an equality law.}
We now compute the finite-rank perturbation in the ideal Gram comparison.
Set $\sigma=0$, take $s=s_*(t)$, and replace each weighted average
$n^{-1}\sum_i A_{ii}f(y_i)$ by
$\E_{N(s,s)}[\sech^2 y\,f(y)]$. Evaluate all other empirical
averages under the same law. Thus $a_y=sb$, $B_A=B=jb$, and
$\kappa^2=1-jq/s$. Write
\[
 X_1=\mathcal Y(\one/\sqrt n),\qquad
 U=\mathcal Y(u),\qquad Z_2=\mathcal Y(Z_1),
\]
using the Gram interpretation just described. Subtracting
$W_0$ from \eqref{eq:conditional-W} gives exactly
\begin{equation}
 \begin{split}
 W-W_0={}&u\bar w^{\mathsf T}+\bar w u^{\mathsf T}
 -(1-\kappa)\{u(W_0u)^{\mathsf T}+(W_0u)u^{\mathsf T}\}\\
 &+\{\alpha+(1-2\kappa)u^{\mathsf T}W_0u\}uu^{\mathsf T}.
 \end{split}
 \label{eq:conditional-rank-difference}
\end{equation}
At equality, $\bar w=j\sqrt{q/s}\,Z_1-Bu$ and $\alpha=2B$ after
discarding the vanishing Gaussian scalar.  Whiten by $S_A$, and use
\eqref{eq:extra-whitened-vector} with $B_A=B$.  The perturbation to
subtract from the identity consists of the planted term
$jX_1X_1^{\mathsf T}$, the term $2jqUU^{\mathsf T}$, and the whitened
version of \eqref{eq:conditional-rank-difference}.  The terms
$-2BUU^{\mathsf T}$ and $+2BUU^{\mathsf T}$ cancel.  Eliminating the
orthogonal direction $N$ by Schur complement adds
$(1-\kappa)^2BUU^{\mathsf T}$ to the perturbation.  Its total
$UU^{\mathsf T}$ coefficient becomes
\[
 2jq-2(1-\kappa)B+(1-\kappa)^2B
 =jq(2-B/s).
\]
Consequently the Schur complement is the identity minus
\begin{equation}
 j\{X_1X_1^{\mathsf T}+(2-B/s)M_1M_1^{\mathsf T}
                 +M_1G_1^{\mathsf T}+G_1M_1^{\mathsf T}\},
 \qquad M_1=\sqrt q\,U,\quad G_1=Z_2/\sqrt s.
 \label{eq:limiting-rank-perturbation}
\end{equation}

To prove positivity of this Schur complement, it suffices to examine
its three-vector Gram matrix. Put
$e=\E_{N(s,s)}[m^2(1-m^2)]$.  The Gram matrix of
$(X_1,M_1,G_1)$ equals
\begin{equation}
 H=\begin{pmatrix}
 b&e&-2e\\
 e&e&b-3e\\
 -2e&b-3e&b/s-2b+6e
 \end{pmatrix}.
 \label{eq:limiting-Gram}
\end{equation}
We derive the entries of this matrix explicitly. The density of
$N(s,s)$ satisfies
$p(y)/p(-y)=e^{2y}$.  Pairing $y$ and $-y$ shows that every integrable
odd function $f$ satisfies $\E f(y)=\E[f(y)\tanh y]$.  This gives
$\E[m(1-m^2)]=e$.  Gaussian integration by parts gives
\[
 \frac1s\E[(y-s)(1-m^2)]=-2e,\qquad
 \frac1s\E[(y-s)m(1-m^2)]=b-3e.
\]
Applying it twice, and using
$(1-m^2)''=4m^2(1-m^2)-2(1-m^2)^2$, gives the final diagonal entry
$b/s-2b+6e$.  The remaining entries are immediate, proving
\eqref{eq:limiting-Gram}.

The coefficient matrix in \eqref{eq:limiting-rank-perturbation} is
\[
 C=j\begin{pmatrix}1&0&0\\0&2-B/s&1\\0&1&0\end{pmatrix}.
\]
Set $d_0=1-jb$ and $g_1=je$.  The first and third columns of $I-HC$
are $(d_0,-g_1,2g_1)^{\mathsf T}$ and
$(-g_1,-g_1,d_0+3g_1)^{\mathsf T}$.  Adding $B/s$ times the third
column to the second makes the second column
$(0,d_0+g_1,0)^{\mathsf T}$.  Hence
\begin{equation}
 \det(I-HC)=(d_0+2g_1)(d_0+g_1)^2>0.
 \label{eq:positive-rank-determinant}
\end{equation}
The same determinant is obtained from the identity minus
\eqref{eq:limiting-rank-perturbation} on any ambient space containing
its three vectors. To determine its sign as a quadratic form, first
subtract only
$jX_1X_1^{\mathsf T}$; the result is positive definite because
$j\norm{X_1}^2=jb<1$.  The remaining two-vector perturbation has at
most one positive eigenvalue, since its coefficient block
$j\left(\begin{smallmatrix}2-B/s&1\\1&0\end{smallmatrix}\right)$
has one positive and one negative eigenvalue. The resulting matrix can
consequently have at most one nonpositive
direction. The positive determinant in
\eqref{eq:positive-rank-determinant} excludes both a zero eigenvalue
and a single negative eigenvalue. Hence the matrix is positive definite.
The Schur complement then also gives positivity before $N$ was eliminated.

\paragraph{Perturbing the scalar law and the diagonal.}
The equality calculation supplies the positive margin we need. We first
explain why the hypotheses put the actual weighted Gram entries near
these equality values.

Each deterministic vector has coordinates $n^{-1/2}$ times
one of $1,m/\sqrt q,(y-d)/\sqrt s$.  Products of two such coordinate
functions are bounded by $C(1+y^2)$.  For any such product $f$, the
tail and Frobenius bounds imply
\begin{equation}
 \left|\frac1n\sum_i(A_{ii}-(V_y)_{ii})f(y_i)\right|
 \le C_{R_1}\frac{\norm{A-V_y}_{\Frob}}{\sqrt n}+C\nu.
 \label{eq:A-moment-comparison}
\end{equation}
On $|y_i|\le R_1$, use Cauchy--Schwarz and boundedness of $f$; outside,
use $|A_{ii}-(V_y)_{ii}|\le2+\eta$ and the tail hypothesis, with
$R_1\ge1$.  The corresponding $V_y$-weighted functions are bounded
and continuous, even when multiplied by $y^2$, because
$\sech^2 y$ decays exponentially.  Their averages are therefore close
to the equality-law averages in a sufficiently small weak
neighborhood.  The parameters $d,s,q,a_y$ are continuous on the
second-moment ball, so their denominators and coefficients are close
as well.  Finally $|B_A-B|\le j\norm{A-V_y}_{\Frob}/\sqrt n$.

We retain the positivity uniformly for $t\in[t_0,T]$ and under
small Gram errors. If $Q$ has finitely many columns, the nonzero
eigenvalues of
$QCQ^{\mathsf T}$ agree with those of
$(Q^{\mathsf T}Q)^{1/2}C(Q^{\mathsf T}Q)^{1/2}$, including
multiplicity.  The largest eigenvalue, with zero adjoined, is therefore
a continuous function of the Gram matrix and the coefficients.
All equality parameters vary continuously on the compact interval
$[t_0,T]$, and $s,q$ stay positive there.  The exact calculation thus
gives a uniform positive whitened margin.  Choose a fixed tolerance
$\xi>0$ on coefficients and Gram entries small enough to preserve,
say, half of it.  Lemmas~\ref{lem:matrix-path-stability}
and~\ref{lem:bilinear}, the Gaussian scalar estimate, and
\eqref{eq:A-moment-comparison} show that this tolerance is achieved
for each fixed admissible $A$ with failure probability at most
$Ce^{-a n}$, after choosing the scalar neighborhood, tail tolerance,
$\sigma_*$, and Frobenius distance sufficiently small.  Returning
from whitening multiplies the margin by at least $c_0$.
The omitted mass $(B-B_A)A$ has norm at most
$j(1+\eta)\norm{A-V_y}_{\Frob}/\sqrt n$ and can be absorbed by
decreasing that distance again.

At this stage the margin is proved for each fixed $A$. To make it
simultaneous, we use that its failure exponent is uniform over the
bounded deterministic vector norms. Those bounds were fixed before
choosing the Frobenius distance.  Also \eqref{eq:conditional-W} implies
$\norm W\le K_1$ with failure $Ce^{-a_1n}$ for some uniform $K_1$:
the norms of the deterministic means are bounded, and the remaining
terms are controlled by $\norm{W_0}$.
On this event, $A\mapsto\mathsf H_{\wt J}(y,A)$ is uniformly
$1/2$-H\"older in diagonal sup norm.  Indeed
$\norm{A^{1/2}-\widehat A^{1/2}}\le
\norm{A-\widehat A}^{1/2}$, while the middle factor in
\eqref{eq:field-hessian} has bounded norm.

Choose a fixed mesh size $\zeta>0$ small enough that the H\"older
error is below one quarter of the margin.  Given $A$ with normalized
Frobenius distance at most $\epsilon$ from $V_y$, let
$\mathcal S_A=\{i:|A_{ii}-(V_y)_{ii}|>\zeta\}$.  It has at most
$\alpha n$ elements, where $\alpha=\epsilon^2/\zeta^2$.
Outside $\mathcal S_A$ replace $A_{ii}$ by $(V_y)_{ii}$; on $\mathcal S_A$, round it
within $\zeta$ to a grid of $[0,1+\eta]$, including the endpoints.
The resulting $\widehat A$ obeys
\[
 \norm{A-\widehat A}\le\zeta,\qquad
 \norm{\widehat A-V_y}_{\Frob}\le2\epsilon\sqrt n.
\]
For fixed $y$, these approximants are deterministic and number at
most
\[
 \sum_{k\le\alpha n}\binom nk N_\zeta^k,
 \qquad N_\zeta\le C(1+1/\zeta).
\]
Its logarithm divided by $n$ tends to zero as $\alpha\downarrow0$;
for instance it is at most
$-\alpha\log\alpha-(1-\alpha)\log(1-\alpha)
 +\alpha\log N_\zeta+o(1)$ when $\alpha<1/2$.
Choose $\epsilon_*$ so small that $2\epsilon_*$ meets all moment
tolerances and this counting rate is less than half the fixed-$A$
failure exponent.  A union bound, followed by the H\"older comparison,
proves the desired uniform margin.

For later use, record the order that preserves the concentration
exponents. First fix $R,t_0,T,\eta$,
the limiting margin, and the Gram tolerance $\xi$.  This fixes the
vector norm bounds, conditional concentration exponents, $K_1$, and
the mesh $\zeta$.  Next take the equality neighborhood and tail
tolerance $\nu$ small.  Choose any sufficiently large finite $R_1$,
then take $\epsilon_*$ small enough for
\eqref{eq:A-moment-comparison} and the union bound.  Finally decrease
$\sigma_*$ so all scalar coefficients lie within the required
tolerance.  No concentration exponent is lost when $\epsilon_*$ is
decreased. This also proves the last assertion of the lemma.
\end{proof}

\subsection{From a small integral to exclusion of approximate zeros}

The conditional estimate controls all nearby diagonals for each fixed
$y$ satisfying the scalar hypotheses. Combined with the scalar rate
bounds, it will make the expected integral over unstable configurations
exponentially small. To exclude approximate zeros at any $y$, the next
deterministic lemma gives the complementary lower bound: an approximate
zero satisfying a regularized determinant bound contributes a definite
exponential amount to the integral over a surrounding ball. The
regularization makes a sign assumption on the Jacobian unnecessary.

\begin{lemma}[Local volume at an approximate zero]
\label{lem:local-field-volume}
Fix $j,K<\infty$, $r>0$, and $\varepsilon>0$.  Suppose
$\norm{\wt J}\le K$, and a set $E\subseteq\R^n$ contains the ball
$\{y':\norm{y'-y}\le r\sqrt n\}$.  Put
\[
 L=F'_{\wt J,h_0}(y)
 =I+(BI-\wt J)V_y-\frac{2j}{n}mm^{\mathsf T}V_y,\qquad
 Q=L^{\mathsf T}L+\gamma I.
\]
If
\begin{equation}
 \frac1n\log\det Q\le jb(y)^2+\varepsilon_1,\qquad
 \norm{F_{\wt J,h_0}(y)}\le\rho\sqrt n,
 \label{eq:local-volume-hypotheses}
\end{equation}
then
\begin{equation}
 (2\pi\sigma^2)^{-n/2}\int_E
 e^{-\norm{F_{\wt J,h_0}(y')}^2/(2\sigma^2)+njb(y')^2/2}\,\dd y'
 \ge e^{-\varepsilon n}
 \label{eq:local-volume-conclusion}
\end{equation}
for all sufficiently large $n$, provided the constants are chosen in
this order: first $\varepsilon_1>0$ sufficiently small in terms of
$\varepsilon$; then any fixed $\gamma>0$; then $\sigma>0$
sufficiently small in terms of $j,K,r,\varepsilon,\gamma$; and finally
$\rho>0$ sufficiently small in terms of those constants and $\sigma$.
The bound is uniform in $y,h_0,n$ under the stated hypotheses.
\end{lemma}

\begin{proof}
We integrate over a Gaussian ellipsoid centered at the approximate
zero. Its covariance is chosen to control the linearized residual.
Let $g''$ be standard Gaussian and put
$\Delta=\sigma Q^{-1/2}g''$, and restrict to
$\mathcal B=\{\norm{g''}\le2\sqrt n\}$.  For
$2\sigma/\sqrt\gamma\le r$, all resulting points $y+\Delta$ belong
to $E$.  Changing variables in the integral gives the lower bound
\begin{equation}
 (\det Q)^{-1/2}\Prob(\mathcal B)
 \E\left[\exp\left\{
 \frac{\norm{g''}^2}{2}
 -\frac{\norm{F(y+\Delta)}^2}{2\sigma^2}
 +\frac{njb(y+\Delta)^2}{2}\right\}\,\middle|\,\mathcal B\right].
 \label{eq:volume-change-variables}
\end{equation}
Here and for the remainder of this proof $F=F_{\wt J,h_0}$.
The conditioning has probability tending to one, and all estimates
below remain valid under it with a uniform factor, since
$\Prob(\mathcal B)\ge1/2$ for large $n$.

We must control the nonlinear residual on this ellipsoid uniformly
in its center. Coordinatewise Taylor expansion gives
\[
 m(y+\Delta)=m+V_y\Delta+r_m,\qquad
 |(r_m)_i|\le C\Delta_i^2.
\]
Writing $b(y+\Delta)=b+\ell_b+r_b$, its linear part is
$\ell_b=-2\av{m(1-m^2)\Delta}$, and
$|r_b|\le C\norm\Delta^2/n$.  Thus, when the product $jb,m$ in
$F$ is expanded, its remainder is bounded in norm by
\[
 C\left(\norm{r_m}+\frac{\norm\Delta^2}{\sqrt n}\right).
\]
For the product of first differences use
$|b(y+\Delta)-b(y)|\le C\norm\Delta/\sqrt n$ and
$\norm{m(y+\Delta)-m(y)}\le\norm\Delta$.  The remaining term
$-\wt Jm$ has remainder bounded by $K\norm{r_m}$.
Since the covariance of $\Delta$ is bounded above by
$\sigma^2I/\gamma$, Gaussian fourth moments give
\[
 \E\sum_i\Delta_i^4\le Cn\sigma^4/\gamma^2,\qquad
 \E\norm\Delta^4\le Cn^2\sigma^4/\gamma^2.
\]
Combining these Taylor estimates with the Gaussian moments, and writing
$R_F=F(y+\Delta)-F(y)-L\Delta$, yields
\begin{equation}
 \E[\norm{R_F}^2\mid\mathcal B]
 \le Cn\sigma^4/\gamma^2,\qquad
 \E[|b(y+\Delta)^2-b(y)^2|\mid\mathcal B]
 \le C\sigma/\sqrt\gamma.
 \label{eq:volume-remainders}
\end{equation}
All constants depend only on $j,K$.

The choice of ellipsoid now controls the linear term: the singular
values of $LQ^{-1/2}$ are at most one. Hence
$\norm{L\Delta}\le\sigma\norm{g''}$ pointwise, including on
$\mathcal B$.  Expanding
$F(y+\Delta)=F(y)+L\Delta+R_F$ and applying Cauchy--Schwarz to its
cross terms yields
\begin{align*}
 &\frac1n\E\left[
 \frac{\norm{g''}^2}{2}
 -\frac{\norm{F(y+\Delta)}^2}{2\sigma^2}
 \,\middle|\,\mathcal B\right]\\
 &\hspace{1cm}\ge
 -C\left\{
 \frac\rho\sigma+\frac\sigma\gamma
 +\left(\frac\rho\sigma\right)^2
 +\left(\frac\sigma\gamma\right)^2
 +\frac\rho\gamma\right\}.
\end{align*}
The last mixed term may also be absorbed into the two squares.
Apply Jensen's inequality to the conditional expectation in
\eqref{eq:volume-change-variables}.  Its logarithm divided by $n$ is
at least
\[
 -\frac{1}{2n}\log\det Q+\frac{jb(y)^2}{2}
 -C\left\{\frac\rho\sigma+\frac\sigma\gamma
          +\left(\frac\rho\sigma\right)^2
          +\left(\frac\sigma\gamma\right)^2
          +\frac\sigma{\sqrt\gamma}\right\}+o(1).
\]
The determinant hypothesis cancels $jb(y)^2/2$, with a loss of at
most $\varepsilon_1/2$. Choose the constants in the order in the
statement; the other errors then have total size smaller than
$\varepsilon$. This proves the lower bound.
\end{proof}

\subsection{Proof of the observation-path statement}

\begin{proof}[Proof of Proposition~\ref{prop:stable-fields}]
We assemble the scalar, matrix, and local-volume estimates under the
planted law \eqref{eq:planted-field}. After excluding unstable
approximate zeros at a fixed time, we remove the auxiliary diagonal,
control the whole path by a fixed mesh, and undo the planting.

Let $J$ be $\wt J$ with its diagonal zeroed. Choose
$\eta,\theta,q_0>0$ sufficiently small that
$j(1+\eta)^2<1$ and, whenever $q\le2q_0$,
\begin{equation}
 I+(j-3jq-2\sqrt j-2\theta)A\succeq c_{\mathrm{sm}}I
 \qquad(0\le A\le(1+\eta)I)
 \label{eq:small-time-matrix-margin}
\end{equation}
for some $c_{\mathrm{sm}}>0$.  These choices are possible because at
$\eta=\theta=q=0$ the lower scalar bound is
$1+j-2\sqrt j=(1-\sqrt j)^2>0$.  On
\begin{equation}
 \norm J\le2\sqrt j+\theta,\qquad
 \max_i|\wt J_{ii}|\le\theta,
 \label{eq:planted-norm-restriction}
\end{equation}
we have $\norm{\wt J}\le2\sqrt j+2\theta$.  Since
$mm^{\mathsf T}/n\preceq qI$, \eqref{eq:small-time-matrix-margin}
implies $\mathsf H_{\wt J}(y,A)\succeq c_{\mathrm{sm}}I$ for all
$q\le2q_0$.  Choose $t_0\in(0,T]$ small enough for
Lemma~\ref{lem:scalar-rates}\textup{(iii)}.

For the remaining times $t\in[t_0,T]$, the conditional Hessian lemma
provides the margin near the equality curve. Choose a second-moment
ball whose complement has a fixed negative rate, using
Lemma~\ref{lem:scalar-rates}\textup{(i)} with $s_{\min}=t_0$.
Apply Lemma~\ref{lem:conditional-field-stability} on this ball.  Take
its equality neighborhood and tail tolerance, choose its tail cutoff
$R_1$ sufficiently large for
Lemma~\ref{lem:scalar-rates}\textup{(iv)}, and then choose its
diagonal-distance radius.  Decrease $\sigma_*>0$ so that
Lemma~\ref{lem:scalar-rates}\textup{(ii)} excludes the complement of
that equality neighborhood, and so that the small-time scalar bound
also holds. The scalar regions excluded by these choices all have fixed
strictly
negative rates, uniformly for $0<\sigma\le\sigma_*$. On the region
left over, the conditional matrix margin fails with probability at
most $e^{-a_*n}$.

Let $\epsilon_*>0$ denote the full diagonal-distance radius just chosen.
Choose $c_E>0$ smaller than both $c_*/2$ and $c_{\mathrm{sm}}/2$.
For $t\ge t_0$, define
\[
 E_t(W)=\left\{y:\exists A,\quad
 0\le A\le(1+\eta)I,\quad
 \norm{A-V_y}_{\Frob}\le\epsilon_*\sqrt n,\quad
 \lambda_{\min}(\mathsf H_{\wt J}(y,A))\le c_E\right\}.
\]
For $0\le t<t_0$, instead put $E_t(W)=\{y:q(y)\ge q_0\}$.
There is $a_0>0$ such that
\begin{equation}
 \E I_{t,\sigma}(E_t)\le e^{-a_0n}
 \quad(0\le t\le T, 0<\sigma\le\sigma_*)
 \label{eq:bad-integral-rate}
\end{equation}
for all sufficiently large $n$, uniformly in $t,\sigma$.
For small times, this follows from the scalar bound and
Lemma~\ref{lem:gaussian-reduction}. For larger times, split the
conditional integral into the excluded scalar regions and their complement.
The excluded regions have negative rates.  On the complement multiply
the conditional probability $e^{-a_*n}$ by the unrestricted scalar
integral, whose rate is at most any prescribed positive number, by
Lemma~\ref{lem:scalar-rates}\textup{(i)}.  Taking that number below
$a_*/2$ proves \eqref{eq:bad-integral-rate}.  Finitely many pieces can
be absorbed by decreasing $a_0$.

We now use local volume to exclude every unstable approximate zero.
An individual witness will force a ball of the excluded set to contribute
too much to the integral. Suppose
\begin{equation}
 \norm{F_{\wt J,h_0}(y)}\le\rho\sqrt n,\qquad
 \norm{A-V_y}_{\Frob}\le\tfrac12\epsilon_*\sqrt n,\qquad
 \lambda_{\min}(\mathsf H_{\wt J}(y,A))\le\tfrac12c_E
 \label{eq:bad-approximate-point}
\end{equation}
for some $0\le A\le(1+\eta)I$, and assume
\eqref{eq:planted-norm-restriction}.  There is a constant $r>0$,
independent of this witness and of $n,t$, such that
$\norm{y'-y}\le r\sqrt n$ implies $y'\in E_t(W)$.
For $t\ge t_0$, the same $A$ witnesses instability throughout a
sufficiently small ball. Indeed,
\[
 \norm{V_{y'}-V_y}_{\Frob}\le2\norm{y'-y},\qquad
 |b(y')-b(y)|\le2\norm{y'-y}/\sqrt n,
\]
and
\[
 \norm{m(y')m(y')^{\mathsf T}/n-m(y)m(y)^{\mathsf T}/n}
 \le2\norm{y'-y}/\sqrt n
\]
preserve the full diagonal-distance constraint and the bound
$\lambda_{\min}\le c_E$ if $r$ is small enough.  For $t<t_0$,
\eqref{eq:small-time-matrix-margin} and the chosen $c_E$ imply
$q(y)>2q_0$; the same Lipschitz bound for $q$ then gives
$q(y')\ge q_0$ for a smaller $r$.

Set $L=F'_{\wt J,h_0}(y)$.  It is a bounded-rank, bounded-norm
perturbation of
\[
 I+(BI-W)V_y:
 \quad L=I+(BI-W)V_y
       -\frac jn\one\one^{\mathsf T}V_y
       -\frac{2j}{n}mm^{\mathsf T}V_y.
\]
Here $n^{-1}\Tr V_y=b(y)$, so the determinant rate in
Lemma~\ref{lem:logdet} matches the weight in \eqref{eq:field-integral}.
That lemma says that, for any chosen $\varepsilon_1>0$, one may choose
a fixed $\gamma>0$ sufficiently
small so that, except on an event of probability $Ce^{-a_1n}$,
simultaneously at every $y$,
\[
 \frac1n\log\det(L^{\mathsf T}L+\gamma I)
 \le jb(y)^2+\varepsilon_1.
\]
The lemma's rank-perturbation allowance applies here with two rank-one
terms of norms at most $j$ and $2j$.
Choose $\varepsilon<a_0/2$, choose $\varepsilon_1$ for
Lemma~\ref{lem:local-field-volume}, then choose $\gamma$ by the
log-determinant lemma.  Next fix $\sigma\in(0,\sigma_*]$ sufficiently
small for the local-volume lemma, and finally fix $\rho>0$
sufficiently small for it.  A witness
\eqref{eq:bad-approximate-point} on the determinant event forces
$I_{t,\sigma}(E_t)\ge e^{-\varepsilon n}$.  Markov's inequality and
\eqref{eq:bad-integral-rate} show that the probability of such a
witness, under \eqref{eq:planted-norm-restriction}, is at most
$Ce^{-a_2n}$, uniformly in $t\in[0,T]$. Here the uniformity of the
negative rates permits the later choice of
$\sigma$ required by the local-volume estimate.

The integral argument used a full GOE matrix. We now pass to its
zero-diagonal version, retaining slack in both residual and Hessian bounds.
Choose a fixed
$\delta_D>0$ so small that
\[
 \delta_D\le\theta,\qquad
 \delta_D\le\rho/2,\qquad
 (1+\eta)\delta_D\le c_E/4.
\]
The diagonal entries of $\wt J$ are independent
$N(j/n,2j/n)$, so
$\Prob(\max_i|\wt J_{ii}|>\delta_D)\le Ce^{-a_Dn}$ for large $n$.
On its complement,
\[
 \frac{\norm{F_{\wt J,h_0}(y)-F_{J,h_0}(y)}}{\sqrt n}
 \le\delta_D,\qquad
 \norm{\mathsf H_{\wt J}(y,A)-\mathsf H_J(y,A)}
 \le(1+\eta)\delta_D.
\]
Thus a zero-diagonal witness with residual at most
$\rho\sqrt n/2$, radius $\epsilon_*\sqrt n/2$, and Hessian minimum
at most $c_E/4$ would give a witness
\eqref{eq:bad-approximate-point}.  We have proved an exponentially
small fixed-time failure probability for the zero-diagonal matrix,
restricted only to $\norm J\le2\sqrt j+\theta$.

To pass from a fixed time to the whole interval, we control observation
increments on a deterministic mesh. Take a mesh of size $\Delta>0$
that includes both endpoints.  For a scalar
Brownian motion, the reflection bound implies that
$\sup_{0\le u\le\Delta}|B(u)|^2/\Delta$ has an exponential moment
at a fixed positive argument.  The coordinate Brownian motions are
independent.  Exponential Markov applied to the sum of their squared
suprema consequently gives, for each desired $\kappa_0>0$, a mesh
$\Delta$ small enough that
\[
 \Prob\left(\max_k\sup_{t_k\le t\le t_{k+1}}
       \frac{\norm{B_0(t)-B_0(t_k)}}{\sqrt n}>\kappa_0\right)
 \le Ce^{-a_3n}.
\]
The deterministic drift contributes at most $\Delta\sqrt n$ to each
increment.  Choose the mesh so the full normalized increment is at
most $\rho/4$.  Set $\rho_0=\rho/8$,
$\epsilon_A=\epsilon_*/2$, and $c=c_E/4$.  A pathwise violation of
\eqref{eq:stable-field} then gives a fixed-time violation at a mesh
point: the Hessian and its witness do not use the field, and the
residual increases by at most $\rho\sqrt n/4$.  A union bound over
the fixed finite mesh proves an annealed path failure bound
$Ce^{-a_4n}$ under the planted law, restricted to
$\norm J\le2\sqrt j+\theta$.

We have obtained a path estimate under planting. To recover the
original disorder law, introduce the partition-function density
\[
 Z(J)=\sum_x e^{x^{\mathsf T}Jx/2},\qquad
 L(J)=Z(J)/\E Z(J).
\]
Size-bias the original joint law of $J,X$, with $X\sim\mu_0$, by
$L(J)$.  Its density in $J,X$ is proportional to
$e^{x^{\mathsf T}Jx/2}$ times the original Gaussian density.
Gaussian tilting makes $X$ uniform and, conditionally on $X$, makes
the off-diagonal entries of $J$ independent with means
$jX_iX_j/n$ and variance $j/n$.  Add independent diagonal entries with law $N(j/n,2j/n)$---a centered
GOE diagonal plus $jI/n$---and conjugate by $D_X$; the resulting matrix and observation have exactly
the law \eqref{eq:planted-field}.  All witness conditions are
equivariant under this conjugation, because $m(D_Xy)=D_Xm(y)$,
$V_{D_Xy}=V_y$, and diagonal $A$ commutes with $D_X$.  Conditional
observation paths given $J$ are unchanged by size bias.

Undoing this change of law requires a lower bound for $L(J)$ on
typical original matrices. For independent uniform signs $\varepsilon_i$,
Gaussian averaging gives
\[
 \frac{\E Z^2}{(\E Z)^2}
 =e^{-j/2}\E\exp\left\{\frac j{2n}
                       \left(\sum_i\varepsilon_i\right)^2\right\}
 \le(1-j)^{-1/2}.
\]
To obtain the inequality, introduce an independent standard Gaussian
$g$ and use
$\E_\varepsilon e^{\sqrt{j/n}g\sum_i\varepsilon_i}
 =\cosh(\sqrt{j/n}g)^n\le e^{jg^2/2}$.
Paley--Zygmund then yields a constant $p_j>0$ with
$\Prob(Z\ge\E Z/2)\ge p_j$.  As a function of the independent
standard off-diagonal Gaussians, $\log Z$ has gradient squared at
most $(j/n)\binom n2\le jn/2$.  Gaussian concentration about a median
therefore has scale $O(\sqrt n)$.  The preceding positive-probability
lower bound implies that this median is at least
$\log\E Z-O_j(\sqrt n)$: otherwise its upper concentration tail
would be smaller than $p_j$.  It follows that, for every fixed
$\delta>0$,
\begin{equation}
 \Prob\{L(J)<e^{-\delta n}\}\longrightarrow0.
 \label{eq:partition-lower-tail}
\end{equation}

We can now turn the annealed planted estimate into the required
conditional observation bound. Let $p_n(J)$ denote the conditional
observation-path failure probability for \eqref{eq:stable-field}.
The planted estimate is
\[
 \E\bigl[L(J)p_n(J)
       \one_{\{\norm J\le2\sqrt j+\theta\}}\bigr]
 \le Ce^{-a_4n}.
\]
Choose $\delta<a_4/2$.  On $L(J)\ge e^{-\delta n}$, this gives an
original expectation bound $Ce^{-(a_4-\delta)n}$ for the restricted
$p_n(J)$.  Markov's inequality makes $p_n(J)\le e^{-an}$ outside an
event of probability tending to zero, for a fixed sufficiently small
$a>0$.  Equation~\eqref{eq:partition-lower-tail} and the usual GOE norm
bound, including diagonal removal, show that the restrictions on
$L(J)$ and $\norm J$ hold with probability tending to one.  Intersect
these events to obtain $\mathcal J_n$ and finish the proof.

We finish by checking measurability of the witness events used in the
argument. At fixed $y$, the matrix-witness sets are projections over
the compact diagonal range.  For existence
of $y,A,t$ along a continuous path, additionally restrict
$\norm y$, $\norm J$, and $\sup_{t\le T}\norm{Y(t)}$ by integers.
The witness inequalities are closed and the witness parameter space
is compact, so these restricted projections are closed.  Taking a
countable union removes the bounds.  This also justifies the
conditional probabilities and the restricted integrals above.
\end{proof}

\subsection{A unique root and quantitative inversion}

The preceding proposition supplies stability along random observations
at fixed subcritical $j$. The following consequence is deterministic:
its hypothesis requires positivity throughout the region of small
residual. From that whole-region assumption we obtain a unique zero
and a bound on the distance to it. This separates the probabilistic
construction of stable fields from the inversion used later.

\begin{lemma}[Root and inverse bound at a stable field]
\label{lem:root}
Fix $j,K,c,r_0>0$.  Let $J$ be symmetric with $\norm J\le K$, and
let $h_0\in\R^n$.  Suppose
\begin{equation}
 \mathsf H_J(y,V_y)\succeq cI
 \quad\text{whenever}\quad
 \norm{F_{J,h_0}(y)}\le r_0\sqrt n.
 \label{eq:root-stability-assumption}
\end{equation}
Then $F_{J,h_0}$ has exactly one zero $r\in\R^n$, and
\begin{equation}
 \norm{y-r}\le C\norm{F_{J,h_0}(y)}
 \quad\text{if}\quad\norm{F_{J,h_0}(y)}<r_0\sqrt n,
 \qquad C=1+(K+3j)/c.
 \label{eq:root-distance}
\end{equation}
In particular this conclusion applies to every field satisfying
\eqref{eq:stable-field}, even if that implication is required only
through residual threshold $\rho_0\sqrt n$.
\end{lemma}

\begin{proof}
We first bound the derivative inverse on the region in the hypothesis.
We then prove existence and uniqueness of a zero, and integrate the
inverse bound along a path to that zero.

Put $L_1=BI-J-2jmm^{\mathsf T}/n$. Since
$F'(y)=I+L_1V_y$, the elementary inverse identity gives
\[
 F'(y)^{-1}
 =I-L_1V_y^{1/2}
       (I+V_y^{1/2}L_1V_y^{1/2})^{-1}V_y^{1/2}.
\]
Under \eqref{eq:root-stability-assumption}, this inverse exists and
has norm at most $1+(K+3j)/c$, since $0<V_y\le I$ and
$\norm{L_1}\le K+3j$.

For existence and uniqueness, pass to the coordinates
$m\in(-1,1)^n$ and define the potential
\[
 \Phi(m)=\sum_i\int_0^{m_i}\atanh z\,\dd z-h_0^{\mathsf T}m
 -\frac12m^{\mathsf T}Jm+\frac j2\norm m^2
 -\frac j{4n}\norm m^4.
\]
Its critical-point equation is
$F_{J,h_0}(\atanh m)=0$.  At any critical point,
\[
 V_y^{1/2}\Phi''(m)V_y^{1/2}
 =\mathsf H_J(y,V_y)\succeq cI,\qquad y=\atanh m,
\]
so the critical point has positive definite Hessian and is isolated.
Choose $a\in(0,1)$ close enough to one that
\[
 \atanh a>1+\max_i\left(|(h_0)_i|+\sum_j|J_{ij}|+j\right).
\]
Every critical point lies in the interior of $[-a,a]^n$, and on each
face of this box the negative gradient $-\nabla\Phi$ points strictly
inward.  The minimum of $\Phi$ on the box is therefore attained at an
interior critical point. This proves existence.

To prove uniqueness, consider the negative gradient flow from any
point of the box. The inward-pointing condition keeps the trajectory
in the box for all positive times.  Its potential decreases and is bounded
below, so
$\int_0^\infty\norm{\nabla\Phi(m(t))}^2\,\dd t<\infty$.
The squared gradient is uniformly continuous in time, because the
gradient and Hessian are bounded on the box.  It follows that the
gradient tends to zero: a nonzero limsup would, by uniform
continuity, force disjoint intervals of fixed positive integral.
Every limit point of a trajectory is therefore critical.
The critical set is finite, since it is a compact set of isolated
zeros.  The limit set of a trajectory is connected: it is the
intersection of the nested compact connected closures of its tails.
Each trajectory therefore converges to one critical point.

It remains to show that all trajectories converge to the same point.
The positive Hessian makes every critical point a strict local
attractor. For clarity, in a sufficiently small ball about it one has
$(m-m_*)^{\mathsf T}\nabla\Phi(m)\ge c_*\norm{m-m_*}^2$,
which gives exponential contraction of the distance along the flow
and invariance of the ball.  Continuous dependence of the flow on
its starting point, up to the time it enters that ball, shows that
each attraction basin is relatively open in $[-a,a]^n$.  These
disjoint nonempty basins cover the connected box.  Hence there is
only one basin and one critical point, proving uniqueness of $r$.

We now derive the distance estimate. Suppose
$\norm{F(y)}<r_0\sqrt n$. Starting at $y(0)=y$, solve on the
open region $\norm{F}<r_0\sqrt n$ the ordinary differential equation
\[
 \dot y(s)=-F'(y(s))^{-1}F(y).
\]
The chain rule gives $F(y(s))=(1-s)F(y)$ while the solution exists.
Its velocity has norm at most $C\norm{F(y)}$.  If its maximal
existence interval ended at or before time one, this velocity bound would
give a finite limiting point, whose residual is still strictly less
than $r_0\sqrt n$ and whose derivative is invertible.  Local existence
would extend the solution, a contradiction. The solution consequently
reaches time one, at which its value is the
unique zero $r$. Integrating the velocity bound along this path gives
\eqref{eq:root-distance}.
\end{proof}
